\scanpage{50}
\mspage{- 3 -}
Ceci étant, il résulte de la définition de $\handwritten{\prod}_{S'/S}$ et de l'existence du
prolongement $u'$ , que $u_{\handwritten{\eta}}$ se prolonge en un $S$-morphisme $u_1:\underset{\handwritten{\scriptscriptstyle\sim\!\sim}}{H}\handwritten{\to}G_1$.
Il reste à voir que $u_1$ se relève en un $S$-morphisme $u:\underset{\handwritten{\scriptscriptstyle\sim\!\sim}}{H}\handwritten{\to}G$ \deleted{$g$} . Comme $\underset{\handwritten{\scriptscriptstyle\sim\!\sim}}{H}$ est connexe, il est loisible de supposer que $p$ est la caractéristique
résiduelle de $S$. Or on a le :

\uline{Lemme 3} : Si $f:G\handwritten{\to}G_1$ désigne le morphisme canonique\inserted{\handwritten{qui induit isom. sur fibres gén.}}, il existe $n\handwritten{\geq}0$,
tel que l'élévation à la puissance $p^n$ dans $G_1$ se factorise en :
\[
G_1\overset{g}{\handwritten{\longrightarrow}}G\overset{f}{\handwritten{\longrightarrow}}G_1.
\]
(ceci a été démontré au cours du séminaire de l'an dernier).

On a alors le diagramme commutatif suivant, dans le quel la première
colonne est exacte :
% Arrows are drawn by hand. Letter and number labels are typed except the outer p^n.
\diagram{\begin{tikzcd}[ampersand replacement=\&,row sep=small,column sep=large]
0\arrow[d] \& \\
H_n\arrow[d] \& \\
\underset{\handwritten{\scriptscriptstyle\sim\!\sim}}{H}\arrow[dd,"p^n"']\arrow[r,"u_1"] \& G_1\arrow[d,"g"]\arrow[dd,bend left=45,"\handwritten{p^n}"] \\
\& G\arrow[d,"f"] \\
\underset{\handwritten{\scriptscriptstyle\sim\!\sim}}{H}\arrow[d]\arrow[ur,dashed,"u"]\arrow[r,"u_1" description] \& G_1\\
\handwritten{0}\&
\end{tikzcd}}

Le morphisme $\deleted{\unsure{f}}\handwritten{g}\,u_1$ s'annule évidemment sur $H_n$ puisqu'il en est ainsi sur
la fibre générique, donc il se factorise à travers $H=H/H_n$ , d'où l'existence du relèvement $u$.

\uline{Remarques} : La démonstration utilise à deux reprises le théorème de Tate ,
qui en toute rigueur n'est démontré que lorsque $\handwritten{\eta}$ est de caractéristique
0 . La première fois , le groupe d'arrivée $G''$ est étale et dans ce cas le
théorème de Tate n'est pas très difficile. Par ailleurs , je sais maintenant
prouver le théorème de \handwritten{T}ate dans le cas où la partie connexe , sur la fibre
générique, dans le groupe d'arrivée, est de type multiplicatif.

\uline{Corollaire} .Soit $G$ un \deleted{$S$}\inserted{$\handwritten{\eta}$}-groupe \deleted{$S$-néronien,} dont la fibre générique est exten%
% archive sheet number: 49
